% Chapter: Conclusions and Future Work
% Continuous Analog Wave Computation in Porous Media

\section{Summary of Contributions}
\label{sec:contributions}

% NOTE: Mirror the objectives from the introduction. State what was actually achieved.
% Be honest: if something is preliminary, say so. Use the same language as the objectives.

[Bullet summary of what was achieved.]

\section{Limitations}
\label{sec:limitations}

% NOTE: Be direct about weaknesses. This strengthens the thesis because it shows critical
% awareness. Do not be apologetic, just factual.

\begin{itemize}
    \item Frequency-domain formulation assumes linear time-invariant steady-state response; transient and nonlinear effects are not captured.
    \item Gradient-free optimisation limits design-space size and convergence speed.
    \item 2D approximation; real porous devices are three-dimensional.
    \item Fixed Lambda and Lambda' reduce the physical realism of the design space.
    \item Mapping from optimised effective properties to a manufacturable geometry is not addressed.
\end{itemize}

\section{Future Work}
\label{sec:future}

% NOTE: Future work should be concrete and extend naturally from the limitations.
% Avoid vague "more research is needed." Specify what kind and why.

\begin{enumerate}
    \item Extend to time-domain FDTD or auxiliary differential equations for broadband transient signals.
    \item Implement adjoint-based gradient optimisation for larger design regions.
    \item Relax fixed Lambda/Lambda' or link them to a pore-scale geometric model.
    \item Validate select designs in OpenFOAM or experimentally.
    \item Demonstrate a concrete classification or routing task with quantified accuracy and generalisation.
\end{enumerate}

\section{Concluding Remarks}
\label{sec:concluding}

% NOTE: One paragraph closing the thesis. Restate the aim, the main finding, and the
% broader implication. End with a forward-looking sentence.

[Final paragraph.]
